\documentclass[10pt,a4paper]{amsart}
\usepackage{amsmath,amssymb,amsthm,graphicx}
% Retained minimal preamble for the Stacks Project chapter more-morphisms.tex (More on Morphisms).
% Verbatim excerpt of the chapter preamble: Xy-pic setup (native preamble lines 14-18) and the native
% theorem declarations with the mathematical macros (native preamble lines 154-217). Copied mechanically;
% nothing added, removed or redefined. The wrapper class and the amsmath/amssymb/amsthm/graphicx base
% packages are supplied by the assembler, so the class/externaldocument part of the native preamble is omitted.
% For commutative diagrams we use Xy-pic
\usepackage[all]{xy}

% We use 2cell for 2-commutative diagrams.
\xyoption{2cell}
% Theorem environments.
%
\theoremstyle{plain}
\newtheorem{theorem}[subsection]{Theorem}
\newtheorem{proposition}[subsection]{Proposition}
\newtheorem{lemma}[subsection]{Lemma}

\theoremstyle{definition}
\newtheorem{definition}[subsection]{Definition}
\newtheorem{example}[subsection]{Example}
\newtheorem{exercise}[subsection]{Exercise}
\newtheorem{situation}[subsection]{Situation}

\theoremstyle{remark}
\newtheorem{remark}[subsection]{Remark}
\newtheorem{remarks}[subsection]{Remarks}

\numberwithin{equation}{subsection}

% Macros
%
\def\lim{\mathop{\mathrm{lim}}\nolimits}
\def\colim{\mathop{\mathrm{colim}}\nolimits}
\def\Spec{\mathop{\mathrm{Spec}}}
\def\Hom{\mathop{\mathrm{Hom}}\nolimits}
\def\Ext{\mathop{\mathrm{Ext}}\nolimits}
\def\SheafHom{\mathop{\mathcal{H}\!\mathit{om}}\nolimits}
\def\SheafExt{\mathop{\mathcal{E}\!\mathit{xt}}\nolimits}
\def\Sch{\mathit{Sch}}
\def\Mor{\mathop{\mathrm{Mor}}\nolimits}
\def\Ob{\mathop{\mathrm{Ob}}\nolimits}
\def\Sh{\mathop{\mathit{Sh}}\nolimits}
\def\NL{\mathop{N\!L}\nolimits}
\def\CH{\mathop{\mathrm{CH}}\nolimits}
\def\proetale{{pro\text{-}\acute{e}tale}}
\def\etale{{\acute{e}tale}}
\def\QCoh{\mathit{QCoh}}
\def\Ker{\mathop{\mathrm{Ker}}}
\def\Im{\mathop{\mathrm{Im}}}
\def\Coker{\mathop{\mathrm{Coker}}}
\def\Coim{\mathop{\mathrm{Coim}}}

% Boxtimes
%
\DeclareMathSymbol{\boxtimes}{\mathbin}{AMSa}{"02}

%
% Macros for moduli stacks/spaces
%
\def\QCohstack{\mathcal{QC}\!\mathit{oh}}
\def\Cohstack{\mathcal{C}\!\mathit{oh}}
\def\Spacesstack{\mathcal{S}\!\mathit{paces}}
\def\Quotfunctor{\mathrm{Quot}}
\def\Hilbfunctor{\mathrm{Hilb}}
\def\Curvesstack{\mathcal{C}\!\mathit{urves}}
\def\Polarizedstack{\mathcal{P}\!\mathit{olarized}}
\def\Complexesstack{\mathcal{C}\!\mathit{omplexes}}
% \Pic is the operator that assigns to X its picard group, usage \Pic(X)
% \Picardstack_{X/B} denotes the Picard stack of X over B
% \Picardfunctor_{X/B} denotes the Picard functor of X over B
\def\Pic{\mathop{\mathrm{Pic}}\nolimits}
\def\Picardstack{\mathcal{P}\!\mathit{ic}}
\def\Picardfunctor{\mathrm{Pic}}
\def\Deformationcategory{\mathcal{D}\!\mathit{ef}}
\usepackage{cleveref}
\crefname{originalreference}{original reference}{original references}
\Crefname{originalreference}{Original reference}{Original references}
\makeatletter
\newlabel{algebra-lemma-refined-quasi-finite-over-polynomial-algebra}{{0520}{source}{}{sourcepointer-1}{}}
\newlabel{algebra-lemma-refined-quasi-finite-over-polynomial-algebra@cref}{{[originalreference][0][]0520}{[1][0][]source}{}{}{}}
\makeatother
\crefname{originallemma}{original lemma}{original lemmas}
\Crefname{originallemma}{Original lemma}{Original lemmas}
\makeatletter
\newlabel{lemma-etale-makes-quasi-finite-finite-at-point}{{02LK}{source}{}{sourcepointer-2}{}}
\newlabel{lemma-etale-makes-quasi-finite-finite-at-point@cref}{{[originallemma][0][]02LK}{[1][0][]source}{}{}{}}
\makeatother
\crefname{originalreference}{original reference}{original references}
\Crefname{originalreference}{Original reference}{Original references}
\makeatletter
\newlabel{morphisms-lemma-finite-presentation-permanence}{{02FV}{source}{}{sourcepointer-3}{}}
\newlabel{morphisms-lemma-finite-presentation-permanence@cref}{{[originalreference][0][]02FV}{[1][0][]source}{}{}{}}
\makeatother
\title{Local structure of a finite type morphism at a point of its fibre}
\author{The Stacks Project Authors}
\date{stacks-project:more-morphisms.tex:local-snapshot-20261011; GNU FDL 1.2 or later}
\begin{document}
\maketitle
\noindent\textit{Editorial note.} Counted span is the author's own proof environment, source lines 13407-13429, immediately following the statement at 13382-13405.\par\medskip
\noindent\textit{Editorial note.} Two cross-chapter pointers (Algebra, Morphisms) plus one same-chapter lemma are recorded in assembly.external\_references.\par\medskip
\noindent\textit{Editorial note.} The statement's four numbered conclusions are all checked inside the span; the closing sentence handles the permanence assertion.\par\medskip
\noindent\textit{Allowed standard prerequisites (editorial).}\begin{itemize}
\item Refined structure of a quasi-finite morphism over a polynomial algebra. Algebra, lemma algebra-lemma-refined-quasi-finite-over-polynomial-algebra, cited in the proof. Produces the morphism to A\textasciicircum{}n\_S with purely transcendental residue field extension that the local model needs.
\item Making quasi-finiteness finite after an elementary etale neighbourhood. chapter lemma lemma-etale-makes-quasi-finite-finite-at-point, cited at proof line 13420. Provides the elementary etale neighbourhood (Y, y) -> (A\textasciicircum{}n\_S, pi(x)) and the open X' of the base change over which the morphism becomes finite.
\item Permanence of finite presentation. Morphisms, lemma morphisms-lemma-finite-presentation-permanence, cited in the proof. Keeps the finite presentation hypothesis along the constructed diagram.
\end{itemize}\medskip
\ifdefined\hypertarget\hypertarget{sourcepointer-1}{}\fi
\noindent\textit{Source reference.} Original reference 0520: focused-seeds/geometry/S9-morph-local-structure-finite-type/ref-tags.txt, source lines 4--4. Navigation pointer only: the counted proof cites this statement by \textbackslash{}ref as a standing input; the deduction the counted proof has to make is carried out inside the counted span, and no part of that deduction is delegated to this pointer.\par\medskip
\ifdefined\hypertarget\hypertarget{sourcepointer-2}{}\fi
\noindent\textit{Source reference.} Original lemma 02LK: sources/stacks-more-morphisms/pinned-original-source.tex, source lines 11924--11924. Navigation pointer only: the counted proof cites this statement by \textbackslash{}ref as a standing input; the deduction the counted proof has to make is carried out inside the counted span, and no part of that deduction is delegated to this pointer.\par\medskip
\ifdefined\hypertarget\hypertarget{sourcepointer-3}{}\fi
\noindent\textit{Source reference.} Original reference 02FV: focused-seeds/geometry/S9-morph-local-structure-finite-type/ref-tags.txt, source lines 5--5. Navigation pointer only: the counted proof cites this statement by \textbackslash{}ref as a standing input; the deduction the counted proof has to make is carried out inside the counted span, and no part of that deduction is delegated to this pointer.\par\medskip
\begin{lemma}
\label{lemma-local-structure-finite-type}
Let $f : X \to S$ be a morphism. Let $x \in X$ and set $s = f(x)$.
Assume that $f$ is locally of finite type and that $n = \dim_x(X_s)$.
Then there exists a commutative diagram
$$
\xymatrix{
X \ar[dd] & X' \ar[l]^g \ar[d]^\pi & x \ar@{|->}[dd] &
x' \ar@{|->}[l]  \ar@{|->}[d] \\
& Y \ar[d]^h & & y \ar@{|->}[d] \\
S \ar@{=}[r] & S & s & s \ar@{=}[l]
}
$$
and a point $x' \in X'$ with $g(x') = x$ such that with $y = \pi(x')$
we have
\begin{enumerate}
\item $h : Y \to S$ is smooth of relative dimension $n$,
\item $g : (X', x') \to (X, x)$ is an elementary \'etale neighbourhood,
\item $\pi$ is finite, and $\pi^{-1}(\{y\}) = \{x'\}$, and
\item $\kappa(y)$ is a purely transcendental extension of $\kappa(s)$.
\end{enumerate}
Moreover, if $f$ is locally of finite presentation then $\pi$ is
of finite presentation.
\end{lemma}

\begin{proof}
The problem is local on $X$ and $S$, hence we may assume that $X$ and
$S$ are affine. By
Algebra, Lemma \ref{algebra-lemma-refined-quasi-finite-over-polynomial-algebra}
after replacing $X$ by a standard open neighbourhood of $x$ in $X$
we may assume there is a factorization
$$
\xymatrix{
X \ar[r]^\pi & \mathbf{A}^n_S \ar[r] & S
}
$$
such that $\pi$ is quasi-finite and such that $\kappa(\pi(x))$
is purely transcendental over $\kappa(s)$. By
Lemma \ref{lemma-etale-makes-quasi-finite-finite-at-point}
there exists an elementary \'etale neighbourhood
$$
(Y, y) \to (\mathbf{A}^n_S, \pi(x))
$$
and an open $X' \subset X \times_{\mathbf{A}^n_S} Y$ which contains a
unique point $x'$ lying over $y$ such that $X' \to Y$ is finite.
This proves (1) -- (4) hold. For the final assertion, use
Morphisms, Lemma \ref{morphisms-lemma-finite-presentation-permanence}.
\end{proof}

\end{document}
